\documentclass[12pt]{article}
\usepackage[utf8]{inputenc}
\usepackage[russian]{babel}
\usepackage{amsmath, amssymb}
\usepackage{geometry}
\usepackage{setspace}
\usepackage{titlesec}

\geometry{a4paper, margin=2cm}
\onehalfspacing
\titleformat{\section}{\large\bfseries\uppercase}{}{0em}{}
\titleformat{\subsection}{\normalsize\bfseries}{}{0em}{}

\title{\textbf{M $\mathbf{\equiv}$ F: Ток как акт тождества между средой и полем}}
\author{ИВин Сергей Юрьевич \\ \texttt{ivin@okean.space}}
\date{2026}

\begin{document}

\maketitle

\begin{abstract}
В данной работе предложено новое основание для понимания электрического тока. Показано, что ток — не движение зарядов, а мгновенный акт тождества между материальной средой и электрическим полем. Утверждается: \( M \equiv F \). Данная модель согласуется с законами Ома и Кирхгофа, переосмысливая их как подтверждение нарушения тождества и его восстановления.
\end{abstract}

\section{Введение}

Классическая физика определяет электрический ток как упорядоченное движение электрических зарядов. Однако данное определение порождает ряд парадоксов:
\begin{itemize}
    \item Почему ток возникает мгновенно по всей цепи?
    \item Почему электроны дрейфуют со скоростью менее 1 мм/с, а «сигнал» распространяется почти со скоростью света?
    \item Где начинается и где заканчивается ток?
\end{itemize}

Все эти вопросы указывают на внутреннюю противоречивость модели движения. В настоящей работе предлагается альтернатива.

\section{Основной принцип: $ M \equiv F $}

Утверждается:

\[
\boxed{M \equiv F}
\]

где:
\begin{itemize}
    \item $ M $ — материальная система (проводник, граница, нагрузка),
    \item $ F $ — электрическое поле на границе этой системы.
\end{itemize}

\textbf{Определение:}  
Ток — это акт, возникающий в момент нарушения тождества $ M \equiv F $ и его последующего восстановления.

Ток не является следствием движения зарядов.  
Ток — это состояние поля на границе, проявляющееся как ответ системы на внешнее воздействие.

\section{Следствия}

\subsection{Закон Ома — как подтверждение нарушения тождества}

Напряжение $ U $ — мера нарушения тождества.  
Сопротивление $ R $ — способность среды восстанавливать $ M \equiv F $.  
Ток $ I = U/R $ — наблюдаемое проявление процесса восстановления.

\subsection{Закон Кирхгофа — как топологическое согласование}

Сумма токов в узле равна нулю — не из-за сохранения заряда,  
а потому что все ветви реагируют синхронно на одно нарушение тождества.

Замкнутый контур — форма поля, в которой восстановление $ M \equiv F $ происходит по цепочке.

\section{Заключение}

Открытие $ M \equiv F $ меняет онтологию тока.  
Ток — не движение,  
не поток,  
не перенос.

Ток — это \textbf{мгновенный акт целостности},  
в котором среда и поле становятся одним.

Данное понимание не отменяет классические законы —  
оно даёт им основание.

\end{document}
